% =============================================================================
% Chapter 16: Responsible AI — ML Systems Lecture Slides (Volume II)
% =============================================================================
% This deck covers responsible AI as engineering infrastructure:
% fairness metrics, impossibility theorems, explainability methods,
% privacy preservation, sociotechnical dynamics, and governance.
%
% IMAGES (SVG → PDF):
%   - responsible-ai-stack.pdf     - fairness-metrics-comparison.pdf
%   - responsible-ai-overhead.pdf  - feedback-loop-bias.pdf
%   - explainability-methods.pdf   - lifecycle-principles.pdf
%   - fairness-tax-visual.pdf
% =============================================================================
\documentclass[aspectratio=169, 12pt]{beamer}
\usepackage{../../assets/beamerthememlsys}

\mlsyssetup{
  volume   = {Volume II},
  chapter  = {Chapter 16},
  logo = {../../assets/img/logo-mlsysbook.png},
  instlogo = {../../assets/img/logo-harvard.png},
  chaptertitle = {Responsible AI},
}

% --- Fonts ---
\usepackage{fontspec}
\setsansfont{Helvetica Neue}[
  BoldFont={Helvetica Neue Bold},
  ItalicFont={Helvetica Neue Italic},
  BoldItalicFont={Helvetica Neue Bold Italic},
]
\setmonofont{JetBrains Mono}[Scale=0.85]

% --- Packages ---
\usepackage{booktabs}
\usepackage{amsmath}

% --- Image paths ---
\graphicspath{
  {images/}
}

% --- Chapter-specific macros ---
\newcommand{\DP}{\text{Demographic Parity}}
\newcommand{\EO}{\text{Equalized Odds}}
\newcommand{\EOpp}{\text{Equality of Opportunity}}

% --- Helper: safe image include ---
\newcommand{\safeimg}[2][width=\textwidth,keepaspectratio]{%
  \IfFileExists{images/#2}{\includegraphics[#1]{#2}}{%
    \IfFileExists{#2}{\includegraphics[#1]{#2}}{%
      \fbox{\parbox[c][2.5cm][c]{0.85\linewidth}{\centering\footnotesize\textcolor{midgray}{[Missing image]}}}%
    }%
  }%
}

% --- Section count for navigation (must match actual \section{} count) ---
\setcounter{mlsystotalsections}{7}

\title{Responsible AI}
\author{Vijay Janapa Reddi}
\institute{Harvard University}
\date{}

\begin{document}

% =============================================================================
% TITLE SLIDE
% =============================================================================
\mlsystitle{Responsible AI}{Ethics as Engineering Infrastructure}{cover_responsible_ai.png}

% =============================================================================
% LEARNING OBJECTIVES
% =============================================================================
\begin{frame}{Learning Objectives}
\note{[2 min] Read through objectives. Emphasize that this chapter transforms ethics
from aspiration into measurable engineering constraints. Ask: ``How many of you have
deployed a model without checking for demographic bias?''}

\footnotesize
\begin{enumerate}\setlength\itemsep{0pt}
  \item Define \textbf{fairness}, \textbf{transparency}, \textbf{accountability}, and \textbf{privacy} as engineering constraints
  \item Calculate \textbf{fairness metrics} (demographic parity, equalized odds, equality of opportunity)
  \item Explain why \textbf{impossibility theorems} prove these metrics are mathematically incompatible
  \item Evaluate \textbf{explainability methods} (SHAP, LIME, gradient-based) and cost-fidelity trade-offs
  \item Analyze \textbf{feedback loops} that amplify bias in production systems
  \item Design \textbf{monitoring infrastructure} for detecting fairness degradation across groups
  \item Assess the \textbf{infrastructure cost} of responsible AI and why it must be budgeted at architecture time
\end{enumerate}

\end{frame}

\begin{frame}{Visual Language}
\note{[1 min] Explain the semantic color system used throughout the course.
These colors are consistent across all diagrams and slides.}

\small
Throughout this course, colors carry meaning:

\vspace{0.3cm}
\begin{columns}[T]
  \begin{column}{0.45\textwidth}
    \begin{mlsyscard}{computestroke}
    \textbf{Blue} --- Compute / Processing\\
    {\footnotesize GPU ops, forward/backward pass, inference}
    \end{mlsyscard}
    \vspace{0.15cm}
    \begin{mlsyscard}{datastroke}
    \textbf{Green} --- Data / Memory\\
    {\footnotesize Data flow, caches, healthy paths}
    \end{mlsyscard}
  \end{column}
  \begin{column}{0.45\textwidth}
    \begin{mlsyscard}{routingstroke}
    \textbf{Orange} --- Routing / Scheduling\\
    {\footnotesize Load balancers, batch windows}
    \end{mlsyscard}
    \vspace{0.15cm}
    \begin{mlsyscard}{errorstroke}
    \textbf{Red} --- Error / Cost / Bottleneck\\
    {\footnotesize Loss, decode phase, waste}
    \end{mlsyscard}
  \end{column}
\end{columns}
\end{frame}



% =============================================================================
\section{The Governance Layer}
% =============================================================================

\begin{frame}{Responsible AI in the Fleet Stack}
\note{[3 min] Position this chapter in the volume. We have built the fleet (Part I),
the distribution logic (Part II), the serving infrastructure (Part III), and the
security armor (earlier in Part IV). Now we give the system a conscience.
Ask: ``What happens when a technically excellent system exhibits bias?''}

% --- Layout: FULL-WIDTH diagram ---
\centering
\safeimg[width=0.88\textwidth,height=4.5cm,keepaspectratio]{responsible-ai-stack.pdf}

\vspace{0.15cm}
\mlsysinsight{Core Thesis:}{Responsibility is infrastructure, not aspiration. Systems that fail responsibility requirements fail to deploy.}

\end{frame}

\begin{frame}{The Amazon Hiring Algorithm}
\note{[3 min] Open with this visceral example. Amazon built a hiring tool trained on
historical resume data. It learned that past successful applicants were predominantly
male and systematically penalized female candidates. It was secure, robust, and
efficient --- yet ethically disastrous.
Ask: ``How would you detect this bias before deployment?''}

\small
\begin{columns}[T]
  \begin{column}{0.55\textwidth}
    \textbf{Amazon (2019):} Hiring algorithm scrapped after
    discovering systematic gender bias.

    \vspace{0.15cm}
    \begin{itemize}\setlength\itemsep{0pt}
      \item Trained on historical resume data
      \item Learned male-dominated hiring patterns
      \item Penalized resumes containing ``women's''
      \item Met every \emph{technical} requirement
    \end{itemize}

    \vspace{0.1cm}
    \begin{mlsyscard}{errorstroke}
    \textbf{Statistically optimal, ethically disastrous.}\\
    {\footnotesize Technical excellence can coexist with social harm.}
    \end{mlsyscard}
  \end{column}
  \begin{column}{0.42\textwidth}
    \vspace{0.3cm}
    \renewcommand{\arraystretch}{1.25}
    {\footnotesize
    \begin{tabular}{@{}ll@{}}
      \toprule
      \textbf{Metric} & \textbf{Status} \\
      \midrule
      Accuracy     & \textcolor{datastroke}{\textbf{Pass}} \\
      Latency      & \textcolor{datastroke}{\textbf{Pass}} \\
      Security     & \textcolor{datastroke}{\textbf{Pass}} \\
      Robustness   & \textcolor{datastroke}{\textbf{Pass}} \\
      \midrule
      Fairness     & \textcolor{errorstroke}{\textbf{FAIL}} \\
      \bottomrule
    \end{tabular}
    }

    \vspace{0.15cm}
    {\footnotesize\textcolor{midgray}{Passing all technical gates\\is necessary but not sufficient.}}
  \end{column}
\end{columns}

\end{frame}

\begin{frame}{From Ethics to Engineering}
\note{[2 min] Key reframing. Traditional engineering treats safety as a guardrail.
Responsible AI is the control plane. An unsafe system is fundamentally unstable:
toxic outputs erode trust (feedback), discrimination degrades training data
(distributional instability), privacy leaks invite regulatory shutdown.}

\footnotesize
\begin{mlsyscard}{crimson}
\textbf{Responsible AI:} Designing, auditing, and operating ML
systems to measurable fairness, safety, privacy, and accountability standards ---
translating ethical principles into verifiable system properties.
\end{mlsyscard}

\vspace{0.1cm}
\renewcommand{\arraystretch}{1.15}
\begin{tabular}{@{}lll@{}}
  \toprule
  & \textbf{Traditional View} & \textbf{Engineering View} \\
  \midrule
  Nature   & Ethical aspiration    & Hard deployment gate \\
  Timing   & Post-hoc review       & Designed-in from data collection \\
  Cost     & ``Nice to have''      & 10--20 ms latency + 15--30\% training \\
  Failure  & Reputational risk     & Regulatory shutdown \\
  \bottomrule
\end{tabular}

\vspace{0.1cm}
\mlsysconcept{Key shift:}{Fairness constraints have the same force as memory limits or latency budgets.}

\end{frame}

% =============================================================================
\section{Core Principles}
% =============================================================================

\begin{frame}{Principles Across the ML Lifecycle}
\note{[3 min] Walk through the lifecycle matrix. Key insight: responsible AI
is a multi-phase commitment. Fairness starts at data collection, not after
deployment. A model trained on biased labels cannot be fixed by post-hoc
threshold adjustment. If short on time, focus on the ``designed-in'' insight.}

% --- Layout: FULL-WIDTH diagram ---
\centering
\safeimg[width=0.92\textwidth,height=4.8cm,keepaspectratio]{lifecycle-principles.pdf}

\vspace{0.1cm}
\mlsysalert{Key:}{Constraints not designed in from data collection require fundamental retraining to fix.}

\end{frame}

\begin{frame}{The Infrastructure Cost of Responsibility}
\note{[3 min] Walk through each row. SHAP is the most expensive: 50--1000x
inference compute. At 10 billion inferences/day, responsible AI overhead can
exceed raw model serving cost. This is not optional --- it is essential
infrastructure like security and redundancy.}

% --- Layout: FULL-WIDTH table diagram ---
\centering
\safeimg[width=0.92\textwidth,height=4.8cm,keepaspectratio]{responsible-ai-overhead.pdf}

\vspace{0.1cm}
{\small Responsible AI is \textbf{provisioned}, not optional --- budgeted like security and fault tolerance.}

\end{frame}

% --- ACTIVE LEARNING 1: Predict ---
\begin{frame}{Predict: What Should ``Fair'' Mean?}
\note{[2 min] Prediction exercise. Give students 60 seconds. The point: there is
no single correct answer. Different definitions conflict mathematically.
This primes the impossibility theorem reveal.}

\centering
\vspace{0.8cm}
{\Large\bfseries Think--Write--Share}

\vspace{0.4cm}
{\large A loan approval model approves Group A at 70\%\\
and Group B at 40\%. Is this unfair?}

\vspace{0.3cm}
{\normalsize Write your definition of ``fair'' in one sentence.}

\vspace{0.3cm}
{\small\textcolor{midgray}{(60 seconds --- then compare with a neighbor.)}}

\vspace{0.2cm}
{\footnotesize\textcolor{lightgray}{Hint: What if Group B has a genuinely higher default rate?}}

\end{frame}

% =============================================================================
\section{Fairness Metrics}
% =============================================================================

\begin{frame}{Three Definitions of Fairness}
\note{[3 min] Reveal after prediction exercise. Walk through each metric.
Demographic Parity: equal outcomes regardless of ground truth.
Equalized Odds: equal error rates conditioned on true label.
Equality of Opportunity: equal TPR among qualified.
The loan example makes the distinctions concrete.}

% --- Layout: FULL-WIDTH comparison ---
\centering
\safeimg[width=0.92\textwidth,height=4.8cm,keepaspectratio]{fairness-metrics-comparison.pdf}

\vspace{0.1cm}
\mlsysalert{Key:}{Each metric captures a \emph{different} aspect of fairness --- and they conflict mathematically.}

\end{frame}

\begin{frame}{Fairness Metrics: The Math}
\note{[3 min] Walk through the formal definitions. Model $h(x)$ predicts binary
outcome, $S$ is a sensitive attribute. Demographic parity ignores ground truth.
Equalized odds conditions on true label. Equality of opportunity focuses on TPR.
Common error: students think one metric implies the others.}

\footnotesize
\textbf{Let} $h(x)$ predict outcome, $S$ = sensitive attribute, $Y$ = true label.

\vspace{0.05cm}
\textbf{1. Demographic Parity} (equal approval rates):
\vspace{-0.1cm}
$$P\big(h(x)=1 \mid S=a\big) = P\big(h(x)=1 \mid S=b\big)$$

\vspace{-0.15cm}
\textbf{2. Equalized Odds} (equal TPR \emph{and} FPR):
\vspace{-0.1cm}
$$P\big(h(x)=1 \mid S=a, Y=y\big) = P\big(h(x)=1 \mid S=b, Y=y\big), \;\; y \in \{0,1\}$$

\vspace{-0.15cm}
\textbf{3. Equality of Opportunity} (equal TPR only):
\vspace{-0.1cm}
$$P\big(h(x)=1 \mid S=a, Y=1\big) = P\big(h(x)=1 \mid S=b, Y=1\big)$$

\vspace{-0.1cm}
\begin{mlsyscard}{errorstroke}
{\scriptsize \textbf{Impossibility} (Kleinberg et al., 2016): When base rates differ, no classifier satisfies all three simultaneously.}
\end{mlsyscard}

\end{frame}

\begin{frame}{The Fairness Tax}
\note{[3 min] Worked example. Unconstrained model: 85\% accuracy, Group A 80\%
approval, Group B 60\% approval. Enforce demographic parity: must approve
Group B at 80\% too. This forces approving risky applicants in Group B.
New accuracy: 81\%. The 4\% drop is the fairness tax.}

% --- Layout: FULL-WIDTH diagram ---
\centering
\safeimg[width=0.88\textwidth,height=4.0cm,keepaspectratio]{fairness-tax-visual.pdf}

\vspace{0.1cm}
\begin{mlsyscard}{crimson}
{\footnotesize Enforcing demographic parity cost \textbf{4\% accuracy} --- a huge drop in credit scoring.
This is the \textbf{Fairness Tax}: the explicit cost of correcting for historical bias.
Engineers must budget this trade-off, not pretend it does not exist.}
\end{mlsyscard}

\end{frame}

% --- ACTIVE LEARNING: Micro-Retrieval Cue ---
\begin{frame}{Quick Check}
\note{[1 min] Answer: No --- impossibility theorem.}

\centering
\vspace{1.0cm}
{\Large\bfseries Quick Check}

\vspace{0.6cm}
{\large Can a classifier satisfy Demographic Parity\\AND Equalized Odds when base rates differ?}

\vspace{0.5cm}
{\normalsize\textcolor{midgray}{15 seconds --- then I will cold-call.}}

\end{frame}



\mlsysfocus{The Impossibility Theorem}{%
When base rates differ between groups,\\[0.3cm]
\textbf{no classifier} can simultaneously satisfy\\[0.2cm]
Calibration $+$ Equalized Odds $+$ Demographic Parity.\\[0.5cm]
{\normalsize Fairness is a \emph{value-laden engineering decision},\\
not a technical optimization with a single correct answer.}%
}

% --- ACTIVE LEARNING 2: Exercise ---
\begin{frame}{Your Turn: Calculate Fairness Metrics}
\note{[3 min] Give students 90 seconds. Group A: 70 approved (40 TP, 30 FP),
30 rejected (5 FN, 25 TN). Group B: 40 approved (30 TP, 10 FP), 60 rejected
(20 FN, 40 TN). Answer: DP gap = 30pp, TPR gap = 29pp, FPR gap = 35pp.}

\footnotesize
\begin{columns}[T]
  \begin{column}{0.48\textwidth}
    {\normalsize\bfseries Calculate Fairness Metrics}

    \vspace{0.15cm}
    \textbf{Group A} (100 applicants):\\
    Approved: 70 (40 repaid, 30 defaulted)\\
    Rejected: 30 (5 would repay, 25 default)

    \vspace{0.1cm}
    \textbf{Group B} (100 applicants):\\
    Approved: 40 (30 repaid, 10 defaulted)\\
    Rejected: 60 (20 would repay, 40 default)

    \vspace{0.1cm}
    Calculate: \textbf{DP gap}, \textbf{TPR gap}, \textbf{FPR gap}.\\
    {\scriptsize\textcolor{midgray}{(90 seconds --- compare with neighbor)}}
  \end{column}
  \begin{column}{0.48\textwidth}
    \pause
    \begin{mlsyscard}{datastroke}
    \textbf{Solution:}\\[0.1cm]
    {\scriptsize
    DP: $70/100 - 40/100$ = \textbf{30 pp}\\[0.05cm]
    TPR: $40/45 - 30/50$ = $0.89 - 0.60$ = \textbf{29 pp}\\[0.05cm]
    FPR: $30/55 - 10/50$ = $0.55 - 0.20$ = \textbf{35 pp}\\[0.1cm]
    \textcolor{errorstroke}{\textbf{Violates all three criteria!}}\\
    Fixing one does not fix the others.
    }
    \end{mlsyscard}
  \end{column}
\end{columns}

\end{frame}

% =============================================================================
\section{Explainability}
% =============================================================================

\begin{frame}{Explainability Methods: Cost vs.\ Fidelity}
\note{[3 min] Walk through the three columns. SHAP: theoretically grounded but
50--1000x compute overhead. LIME: model-agnostic local explanations, 100--1000
forward passes. Gradient-based: cheapest but only works on differentiable models
and lowest fidelity. Selection depends on deployment constraints.}

% --- Layout: FULL-WIDTH diagram ---
\centering
\safeimg[width=0.92\textwidth,height=4.8cm,keepaspectratio]{explainability-methods.pdf}

\vspace{0.1cm}
\mlsysconcept{Selection rule:}{Match method to deployment constraint --- GDPR requires SHAP; real-time needs gradients.}

\end{frame}

\begin{frame}{The Explainability Budget}
\note{[2 min] Concrete worked example. A recommendation system at 100ms SLA,
10{,}000 QPS. SHAP adds 50--200ms per request. Total: 150--300ms. SLA violated.
Options: pre-compute explanations, use gradient methods, or accept 5x infra cost.
Key: explainability must be budgeted at architecture time.}

\footnotesize
\begin{columns}[T]
  \begin{column}{0.55\textwidth}
    \textbf{Scenario:} Rec.\ system, SLA: 100 ms, 10{,}000 QPS

    \vspace{0.1cm}
    \begin{mlsyscard}{errorstroke}
    {\scriptsize
    \textbf{SHAP adds}: 50--200 ms/request\\
    \textbf{Total latency}: 150--300 ms --- \textcolor{errorstroke}{\textbf{SLA VIOLATED}}}
    \end{mlsyscard}

    \vspace{0.1cm}
    \textbf{Options:}
    \begin{itemize}\setlength\itemsep{0pt}
      \item Pre-compute explanations (offline)
      \item Use gradient-based methods (+1--10$\times$)
      \item Accept 5$\times$ infrastructure cost
      \item Choose inherently interpretable model
    \end{itemize}
  \end{column}
  \begin{column}{0.42\textwidth}
    \renewcommand{\arraystretch}{1.1}
    {\scriptsize
    \begin{tabular}{@{}lr@{}}
      \toprule
      \textbf{Method} & \textbf{Overhead} \\
      \midrule
      SHAP            & 50--1000$\times$ \\
      LIME            & 100--1000 passes \\
      Gradient        & 1--10$\times$ \\
      Inherent        & 0$\times$ (arch. cost) \\
      \bottomrule
    \end{tabular}
    }

    \vspace{0.1cm}
    {\scriptsize\textcolor{midgray}{GDPR Article 22: Automated decisions\\
    with legal effects \emph{must} be explained.\\
    Fines: \textbf{\texteuro 4.5B+} cumulative.}}
  \end{column}
\end{columns}

\end{frame}

% =============================================================================
\section{Feedback Loops}
% =============================================================================

\begin{frame}{Feedback Loops Amplify Bias}
\note{[3 min] The predictive policing example. Historical data over-represents
certain neighborhoods. Model predicts ``high crime'' there. More patrols sent.
More arrests made. Training data reinforced. The system destabilizes its own
data distribution. Monitoring closes the governance loop.}

% --- Layout: FULL-WIDTH comparison ---
\centering
\safeimg[width=0.92\textwidth,height=4.8cm,keepaspectratio]{feedback-loop-bias.pdf}

\vspace{0.1cm}
\mlsysalert{Unstable:}{An unfair system degrades its own training data --- bias is a distributional instability.}

\end{frame}

\begin{frame}{Sociotechnical Dynamics}
\note{[2 min] Four challenges beyond pure technology. Automation bias: humans
over-trust the model. Value conflicts: accuracy vs.\ fairness vs.\ privacy.
Contestability: affected individuals must challenge decisions. Institutional
embedding: governance structures must outlast individual champions.}

\small
\begin{columns}[T]
  \begin{column}{0.48\textwidth}
    \begin{mlsyscard}{computestroke}
    \textbf{Automation Bias}\\
    {\footnotesize Humans over-trust model outputs. Radiologists shown AI predictions exhibit \textbf{confirmation bias}, accepting errors they would catch independently.}
    \end{mlsyscard}

    \vspace{0.15cm}

    \begin{mlsyscard}{routingstroke}
    \textbf{Value Conflicts}\\
    {\footnotesize Accuracy vs.\ fairness vs.\ privacy are competing constraints. No single optimization satisfies all stakeholders.}
    \end{mlsyscard}
  \end{column}
  \begin{column}{0.48\textwidth}
    \begin{mlsyscard}{datastroke}
    \textbf{Contestability}\\
    {\footnotesize Affected individuals must have mechanisms to challenge automated decisions. Requires explanation infrastructure.}
    \end{mlsyscard}

    \vspace{0.15cm}

    \begin{mlsyscard}{errorstroke}
    \textbf{Institutional Embedding}\\
    {\footnotesize Governance structures must outlast individual champions. Responsible AI fails when tied to a single team.}
    \end{mlsyscard}
  \end{column}
\end{columns}

\end{frame}

% --- ACTIVE LEARNING 3: Discussion ---
\begin{frame}{Discussion: Which Fairness Metric?}
\note{[3 min] Turn-and-talk. Students discuss in pairs. Cold-call 2--3 pairs.
There is no single right answer. Criminal justice favors equalized odds
(equal error rates). Hiring favors demographic parity (four-fifths rule).
Healthcare favors equality of opportunity (equal treatment for equally sick).}

\centering
\vspace{0.6cm}
{\large\bfseries Turn and Talk \textcolor{midgray}{(90 seconds)}}

\vspace{0.4cm}
{\large You are deploying a model in \textbf{three} domains:\\[0.15cm]
\textbf{1.} Criminal justice risk scoring\\
\textbf{2.} Employee hiring screening\\
\textbf{3.} Healthcare triage\\[0.2cm]
\alert{Which fairness metric would you prioritize for each?}}

\vspace{0.3cm}
{\small\textcolor{midgray}{Consider: What type of error causes the most harm?}}

\end{frame}

% =============================================================================
\section{Privacy \& Governance}
% =============================================================================

\begin{frame}{Privacy Preservation at Scale}
\note{[3 min] Two key techniques. Differential privacy: mathematical guarantee
that any individual's data has bounded influence. Cost: 2--5\% accuracy loss,
15--30\% training overhead. Machine unlearning: remove specific user influence
from trained models without retraining from scratch. GDPR right to be forgotten.}

\footnotesize
\begin{columns}[T]
  \begin{column}{0.48\textwidth}
    \begin{mlsyscard}{routingstroke}
    \textbf{Differential Privacy}\\[0.05cm]
    {\scriptsize
    Mathematical guarantee: individual's data has \textbf{bounded influence} on model output.\\[0.05cm]
    \textbf{DP-SGD}: Clip gradients + add noise.\\[0.05cm]
    \textbf{Cost}: 2--5\% accuracy, +15--30\% training.\\[0.05cm]
    \textbf{Benefit}: Enables legally restricted data.}
    \end{mlsyscard}
  \end{column}
  \begin{column}{0.48\textwidth}
    \begin{mlsyscard}{computestroke}
    \textbf{Machine Unlearning}\\[0.05cm]
    {\scriptsize
    Remove specific user's influence from a trained model.\\[0.05cm]
    \textbf{Exact}: Retrain from scratch (expensive).\\[0.05cm]
    \textbf{Approximate}: Influence functions.\\[0.05cm]
    \textbf{Legal}: GDPR ``right to be forgotten,'' CCPA.}
    \end{mlsyscard}
  \end{column}
\end{columns}

\vspace{0.05cm}
\mlsysconcept{Key:}{Privacy is a temporal dimension --- unlearning enables compliance without full retraining.}

\end{frame}

\begin{frame}{Governance Infrastructure}
\note{[2 min] Four pillars of governance. Model cards document intended use,
limitations, and performance across groups. Audit trails enable traceability.
System prompts (for LLMs) are the first line of defense. All require CI/CD
rigor equivalent to model weights.}

\scriptsize
\renewcommand{\arraystretch}{1.1}
\begin{tabular}{@{}p{2.2cm}p{3.8cm}p{4cm}@{}}
  \toprule
  \textbf{Mechanism} & \textbf{Purpose} & \textbf{Systems Requirement} \\
  \midrule
  \textbf{Model Cards}     & Document use, limits, group perf. & Version-controlled docs in CI/CD \\
  \textbf{Audit Trails}    & Decision and data lineage tracing & +2--5 ms/inference, TB-scale logs \\
  \textbf{System Prompts}  & First line of LLM behavior defense & Version control across M users \\
  \textbf{Red Teaming}     & Probe for bias, toxicity, jailbreaks & Specialized adversarial teams \\
  \bottomrule
\end{tabular}

\vspace{0.1cm}
\begin{mlsyscard}{crimson}
{\scriptsize Managing system prompts across a fleet requires the same \textbf{version control and CI/CD rigor} as model weights. Governance is infrastructure.}
\end{mlsyscard}

\end{frame}

% =============================================================================
\section{Implementation}
% =============================================================================

\begin{frame}{Implementation Challenges}
\note{[2 min] Five barriers to implementing responsible AI in practice.
Organizational: ethics teams lack authority. Data: protected attributes
often missing. Competing objectives: accuracy vs.\ fairness vs.\ cost.
Scalability: monitoring overhead compounds at fleet scale.
Standardization: no agreed-upon benchmarks for fairness evaluation.}

\footnotesize
\begin{columns}[T]
  \begin{column}{0.48\textwidth}
    \textbf{Organizational Barriers:}
    \begin{itemize}\setlength\itemsep{-1pt}
      \item Ethics teams lack deployment authority
      \item Responsibility tied to individual champions
      \item Incentive structures misaligned
    \end{itemize}

    \vspace{0.05cm}
    \textbf{Data Constraints:}
    \begin{itemize}\setlength\itemsep{-1pt}
      \item Protected attributes often missing
      \item Ground truth labels reflect historical bias
      \item Proxy variables encode structural inequity
    \end{itemize}
  \end{column}
  \begin{column}{0.48\textwidth}
    \textbf{Competing Objectives:}
    \begin{itemize}\setlength\itemsep{-1pt}
      \item Accuracy vs.\ fairness (4\% tax)
      \item Privacy vs.\ transparency
      \item Speed-to-market vs.\ governance
    \end{itemize}

    \vspace{0.05cm}
    \textbf{Scale \& Standards:}
    \begin{itemize}\setlength\itemsep{-1pt}
      \item Monitoring overhead compounds at fleet scale
      \item No agreed fairness benchmarks
      \item Regulatory fragmentation (GDPR vs.\ CCPA)
    \end{itemize}
  \end{column}
\end{columns}

\end{frame}

\begin{frame}{AI Safety and Value Alignment}
\note{[2 min] The frontier. Value alignment: ensuring AI optimizes for human values,
not proxy objectives. YouTube pre-2017: optimized for CTR, promoted conspiracy content.
RLHF bridges human preference to model weights. System prompts control LLM behavior.
These are the primary sociotechnical mechanisms for the generative era.}

\footnotesize
\begin{columns}[T]
  \begin{column}{0.55\textwidth}
    \textbf{Value Alignment Problem:}\\
    {\scriptsize AI systems may optimize proxy objectives diverging from human intent.}

    \vspace{0.1cm}
    \begin{mlsyscard}{errorstroke}
    {\scriptsize \textbf{YouTube (pre-2017):} Optimized for CTR $\to$ promoted conspiracy content that maximized clicks while degrading user welfare. Required full reward pipeline redesign.}
    \end{mlsyscard}

    \vspace{0.1cm}
    \textbf{Generative Alignment:}\\
    {\scriptsize RLHF bridges human preferences to model weights. System prompts govern fleet behavior.}
  \end{column}
  \begin{column}{0.42\textwidth}
    \renewcommand{\arraystretch}{1.1}
    {\scriptsize
    \begin{tabular}{@{}lr@{}}
      \toprule
      \textbf{Mechanism} & \textbf{Scope} \\
      \midrule
      RLHF          & Training time \\
      System prompts & Serving time \\
      Red teaming    & Pre-deploy \\
      Monitoring     & Continuous \\
      \bottomrule
    \end{tabular}
    }

    \vspace{0.1cm}
    {\scriptsize\textcolor{midgray}{RLHF is limited by the representativeness of the rating population.}}
  \end{column}
\end{columns}

\end{frame}

% =============================================================================
\section{Wrap-Up}
% =============================================================================

\begin{frame}{Fallacies}
\note{[2 min] Four misconceptions with quantitative evidence.}

\footnotesize
\textbf{Fallacy:} \textit{Bias can be eliminated through better algorithms and more data.}\\
{\scriptsize The healthcare algorithm affected 200M Americans; the issue was proxy selection, not data quantity. Impossibility theorems prove no algorithm satisfies all fairness criteria simultaneously.}

\vspace{0.08cm}
\textbf{Fallacy:} \textit{Achieving one fairness metric guarantees overall fairness.}\\
{\scriptsize Demographic parity forces calibration violations. Single-metric optimization exposes organizations to litigation on other legally relevant criteria.}

\vspace{0.08cm}
\textbf{Fallacy:} \textit{Responsible AI practices impose only costs without business value.}\\
{\scriptsize Differential privacy enables use of legally restricted data. Bias monitoring detected the healthcare algorithm's 50\% enrollment reduction before regulatory intervention.}

\end{frame}

\begin{frame}{Pitfalls}
\note{[2 min] Three operational pitfalls.}

\footnotesize
\textbf{Pitfall:} \textit{Treating explainability as an optional post-deployment feature.}\\
SHAP adds 50--200 ms per request. A 100 ms SLA system cannot retrofit SHAP without 5$\times$ infrastructure cost. Explainability must be budgeted at architecture time.

\vspace{0.15cm}
\textbf{Pitfall:} \textit{Implementing fairness via post-hoc threshold adjustment.}\\
Group-specific thresholds destroy calibration: ``80\% confidence'' means different things for different groups. Proper approach: jointly optimize during training, accepting accuracy-fairness trade-offs.

\vspace{0.15cm}
\textbf{Pitfall:} \textit{Assuming governance structures survive without institutional embedding.}\\
Responsible AI tied to individual champions fails when they leave. Governance must be embedded in CI/CD pipelines, deployment gates, and organizational incentive structures.

\end{frame}

% --- RETRIEVAL PRACTICE ---

% --- MUDDIEST POINT ---
\begin{frame}{Muddiest Point}
\note{[2 min] Quick anonymous poll. Students write on a slip of paper or submit
digitally. Collect and scan for patterns. Address the top 2--3 confusions in the
next lecture's opening. This closes the feedback loop.}

\centering
\vspace{1.0cm}
{\Large\bfseries What was the \alert{muddiest point} today?}

\vspace{0.8cm}
{\normalsize Write down the concept you found \textbf{most confusing}.}

\vspace{0.5cm}
{\small\textcolor{midgray}{Anonymous. One sentence. Submit before you leave.}}

\end{frame}

\begin{frame}{What Were the Key Ideas?}
\note{[2 min] Retrieval practice. Students write 90 seconds, no notes.
Walk around the room.}

\centering
\vspace{1.5cm}
{\Large\bfseries Close your notes.}

\vspace{0.8cm}
{\large Write down the \textbf{4 most important concepts} from today.}

\vspace{0.8cm}
{\normalsize\textcolor{midgray}{90 seconds --- no peeking.}}

\end{frame}

\begin{frame}{Key Takeaways}
\note{[2 min] Reveal. Walk through each bullet. Emphasize: impossibility theorem,
fairness tax, infrastructure cost, feedback loops, governance as CI/CD.}

\scriptsize
\begin{itemize}\setlength\itemsep{0pt}
  \item \textbf{Responsibility is infrastructure}: Fairness monitoring adds 10--20 ms; SHAP costs 50--1000$\times$ inference compute. Budget at architecture time.
  \item \textbf{Impossibility theorem}: No classifier satisfies Calibration + Equalized Odds + Demographic Parity when base rates differ.
  \item \textbf{The Fairness Tax}: Enforcing demographic parity costs $\sim$4\% accuracy --- the explicit price of correcting historical bias.
  \item \textbf{Feedback loops}: Biased systems destabilize their own training data. Monitoring closes the governance loop.
  \item \textbf{Explainability spectrum}: SHAP (high fidelity, 50--1000$\times$) to gradient-based (low cost, lower fidelity).
  \item \textbf{Privacy preservation}: Differential privacy (2--5\% accuracy, +15--30\% training) enables legally restricted data use.
  \item \textbf{Governance as CI/CD}: Model cards, audit trails, system prompts need the same version control rigor as model weights.
\end{itemize}

\end{frame}

\begin{frame}{References}
\note{[1 min] Point students to canonical papers.}

\small
\mlsysref{Kleinberg+16}{Kleinberg et al. ``Inherent Trade-Offs in the Fair Determination of Risk Scores.'' 2016.}
\mlsysref{Chouldechova17}{A. Chouldechova. ``Fair Prediction with Disparate Impact.'' 2017.}
\mlsysref{Obermeyer+19}{Obermeyer et al. ``Dissecting Racial Bias in Healthcare.'' \textit{Science}, 2019.}
\mlsysref{Hardt+16}{Hardt, Price, Srebro. ``Equality of Opportunity in Supervised Learning.'' NeurIPS, 2016.}
\mlsysref{Ribeiro+16}{Ribeiro, Singh, Guestrin. ```Why Should I Trust You?' LIME.'' KDD, 2016.}
\mlsysref{Lundberg+17}{Lundberg \& Lee. ``A Unified Approach to Interpreting Model Predictions.'' NeurIPS, 2017.}

\end{frame}

\begin{frame}{Next Lecture: Conclusion}
\note{[1 min] Forward hook. We have built the entire stack: fleet, distribution,
deployment, and governance. The conclusion synthesizes all principles into
a unified engineering framework for distributed ML systems.}

\centering
\vspace{0.5cm}

{\large We have addressed every dimension of production ML systems:}

\vspace{0.4cm}
\begin{columns}[c]
  \begin{column}{0.22\textwidth}\centering\small
    \textcolor{computestroke}{\textbf{Performance}}\\[0.1cm]
    {\footnotesize Parts I--III}
  \end{column}
  \begin{column}{0.22\textwidth}\centering\small
    \textcolor{routingstroke}{\textbf{Security}}\\[0.1cm]
    {\footnotesize Ch 13}
  \end{column}
  \begin{column}{0.22\textwidth}\centering\small
    \textcolor{datastroke}{\textbf{Sustainability}}\\[0.1cm]
    {\footnotesize Ch 15}
  \end{column}
  \begin{column}{0.22\textwidth}\centering\small
    \textcolor{errorstroke}{\textbf{Responsibility}}\\[0.1cm]
    {\footnotesize Ch 16 (today)}
  \end{column}
\end{columns}

\vspace{0.5cm}
{\normalsize In \textbf{Ch 17: Conclusion}, we synthesize the principles\\
from across this volume into a \textbf{unified perspective}\\
on distributed ML systems engineering.}

\vspace{0.3cm}
{\small\textbf{The enduring lessons that guide practice regardless of technology.}}

\end{frame}



\appendix

\begin{frame}{Backup: Extended Reference}
\note{Backup slide with additional reference material for this chapter.}

\footnotesize
This slide provides extended reference material for students who want to go deeper.
Refer to the textbook chapter for complete derivations and additional worked examples.

\vspace{0.3cm}
\begin{mlsyscard}{crimson}
\textbf{Key equations and frameworks from this chapter} are collected in the
textbook's summary tables. Use them as a quick reference during problem sets.
\end{mlsyscard}

\end{frame}

\begin{frame}{Backup: Further Reading}
\note{Backup slide. Point students to additional resources beyond the references slide.}

\footnotesize
\textbf{For deeper exploration:}

\vspace{0.15cm}
\begin{itemize}\setlength\itemsep{2pt}
  \item The textbook chapter contains 2--3 additional worked examples
  \item Problem sets apply these concepts to real-world scenarios
  \item The course website links to open-source implementations
  \item Office hours: bring your toughest questions
\end{itemize}

\end{frame}


\end{document}
